\documentclass{article}
\usepackage[T1]{fontenc}
\usepackage{lmodern}
\usepackage{graphicx}
\usepackage{amsmath,amssymb}
\usepackage[a4paper,margin=2cm]{geometry}
\usepackage[absolute,overlay]{textpos}
\setlength{\TPHorizModule}{1mm}
\setlength{\TPVertModule}{1mm}
\usepackage[indonesian]{babel}
\usepackage{hyperref}
\setlength{\emergencystretch}{2em}
% unit: Halaman 36

\begin{document}

% === Halaman 36 (llm) ===

\begin{itemize}
    \item \textbf{XOR Rcon}: XOR-kan hasilnya dengan \textit{round constant} (Rcon) ke $i/4$ (atau $Rcon[i/4]$).
    \begin{itemize}
        \item Nilai $Rcon$ berbeda-beda untuk setiap $j = i/4$, yaitu $Rcon[j] = (RC[j], 0, 0, 0)$, dengan $RC[1]=1, RC[j] = 2 \bullet RC[j-1]$, simbol $\bullet$ menyatakan perkalian yang didefinisikan di dalam $\text{GF}(2^8)$.
        \item Nilai $RC[j]$ di dalam heksadesimal adalah [STA11]: $RC[1]=01, RC[2]=02, RC[3]=04, RC[4]=08, RC[5]=10, RC[6]=20, RC[7]=40, RC[8]=80, RC[9]=1B, RC[10]=36$.
    \end{itemize}
    \item Simpan hasil fungsi $g$ ke dalam peubah \textit{temp}
\end{itemize}

\begin{enumerate}[label=\alph*)]
    \setcounter{enumi}{2}
    \item XOR-kan $w[i-4]$ dengan \textit{temp}
\end{enumerate}

\end{document}
